\documentclass[a4paper,12pt]{article}
\usepackage{graphicx}
\usepackage{caption}
\usepackage{subcaption}
\usepackage{float}
\usepackage[margin=1in]{geometry}
\usepackage{lmodern}
\usepackage{amsmath}

\begin{document}

\title{Relatório RBF}
\author{Eduardo}
\date{\today}
\maketitle

\section*{Introduction}
This report presents the implementation and evaluation of a Radial Basis Function (RBF) classifier with L2 regularization. The model leverages K-means clustering for center selection and employs regularized least squares for weight optimization. Experiments were conducted on synthetic and real-world datasets to validate performance.

\section*{Methodology}
\subsection*{RBF Network Structure}
The classifier computes activations using the RBF kernel:
\begin{equation}
  \phi(\mathbf{x}) = \exp(-\gamma \|\mathbf{x} - \mathbf{c}_i\|^2)
\end{equation}
where \(\mathbf{c}_i\) are cluster centers and \(\gamma\) is determined by:
\begin{equation}
  \gamma = \frac{1}{2 \cdot (\text{median}(\text{min\_dists}))^2}
\end{equation}
with fallback to data-to-center distances if centers collapse.

\subsection*{Weight Optimization}
Weights \(\mathbf{W}\) are solved using regularized least squares:
\begin{equation}
  \mathbf{W} = (\mathbf{H}^T\mathbf{H} + \lambda \mathbf{I})^{-1}\mathbf{H}^T\mathbf{Y}
\end{equation}
where \(\mathbf{H}\) is the RBF activation matrix and \(\lambda\) is the regularization parameter.

\subsection*{Implementation Details}
\begin{itemize}
  \item Centers selected via K-means clustering
  \item Automatic \(\gamma\) computation using median heuristic
  \item L2 regularization for numerical stability
  \item One-vs-rest strategy for multi-class support
\end{itemize}

\section*{Experiments}
\subsection*{Synthetic Data}
\begin{itemize}
  \item 1,000 samples, 20\% test split
  \item Default parameters: 100 centers, \(\lambda=10^{-3}\)
\end{itemize}

\subsection*{Breast Cancer Dataset}
\begin{itemize}
  \item Preprocessing: Drop correlated features
  \item 10-fold cross-validation
  \item Evaluation metric: Classification accuracy
\end{itemize}

\section*{Results}
\subsection*{Synthetic Data Evaluation}
\begin{equation*}
  \text{Accuracy} = 0.88
\end{equation*}

\subsection*{Cross-Validation Performance}
\begin{table}[H]
  \centering
  \begin{tabular}{|l|c|}
  \hline
  Mean Accuracy & 0.89 \\
  Std Deviation & 0.03 \\
  \hline
  \end{tabular}
  \caption{10-fold cross-validation results on breast cancer dataset}
\end{table}

Detailed fold scores: [0.91, 0.89, 0.88, 0.95, 0.86, 0.88, 0.84, 0.88, 0.88, 0.95]

\section*{Conclusion}
The RBF classifier demonstrated strong performance on both synthetic (88\% accuracy) and real-world data (89\% mean cross-validated accuracy). The median-based \(\gamma\) computation and L2 regularization proved effective for stable learning.

\end{document}